% SOURCE_AUTHORITY: sga5_fr_workpass.tex, lines 12919--12925 (LF-normalized unit).
% SOURCE_SHA256: 791F4EFFC5E02832D5D77ED03518C8156D6F07E4C8238B03545DB93D883FBB28
Demos agora exemplos de categorias $\Delta$ e de complexos de feixes $F$ para os quais se verificam as condições enunciadas anteriormente.

a) Sejam $\Delta=D^b(A)_{\coh}$ (VIII 5) e $F$ um complexo de feixes de $A$-módulos construtível (SGA 4 XVII) e de cohomologia limitada. Neste caso, $\varepsilon_x^\Delta(F)\in K_\bullet(A)$.

b) Toma-se $\Delta=D(A)_{\parf}$ (VIII 6) e $F$ um complexo de feixes de $A$-módulos construtível e de dimensão Tor finita. Neste caso, $\varepsilon_x^\Delta(F)\in K^{\bullet}(A)$.

c) Sejam $A$, $B$ duas $\Lambda$-álgebras e $\varphi:A\to B$ um homomorfismo de $\Lambda$-álgebras. Toma-se $\Delta=D^b(B\operatorname{ rel }A)_{\coh}$ (VIII 10) e $F$ um complexo de feixes de $B$-módulos, construtível como feixe de complexos de $A$-módulos e de cohomologia limitada. Obtém-se $\varepsilon_x^\Delta(F)\in K_\bullet(B\operatorname{ rel }A)$. Evidentemente, pode-se variar como antes para obter, por exemplo, elementos $\varepsilon_x^\Delta(F)\in K^{\bullet}(B\operatorname{ rel }A)$.
